% GENERATED FILE. Edit canonical lessons, then run npm run generate:books.
% canonical-source-hash: fnv1a64:fe5d276065888ce5
% canonical-lessons: AR-C543-lahn, AR-C543-iqa, AR-C543-azafa, AR-C543-mulahhin, AR-C543-jawqa, AR-R543-review-536-543

\chapter{Making music}
\label{ch:ar-making-music}
\hlchaptermodality{543}

\begin{chapteropening}
\textbf{By the end of this chapter:} \emph{I can talk about a melody, a rhythm, a composer and a choir, and say that someone played an instrument.}

Complete the last lesson of chapter 543: I can talk about a melody, a rhythm, a composer and a choir, and say that someone played an instrument.
\end{chapteropening}

\section[\texorpdfstring{\ar{لحن} (laḥn)}{لحن (laḥn)}]{\ar{لحن} (laḥn) — a melody, a tune}

\label{lesson:AR-C543-lahn}



\begin{quote}
Before the new one: say the Arabic for to applaud, then the Arabic for acting.
\end{quote}

\subsection*{You'll want to know: \ar{لحن}}
\textbf{\ar{لحن}} — \emph{laḥn} — \textquotedblleft{}a melody, a tune\textquotedblright{}.

\begin{quote}
\ar{غنت} \ar{البنت} \ar{لحنا} \ar{جميلا}.
\end{quote}

\begin{quote}
\emph{ghannat al-bint laḥnan jamīlan.}
\end{quote}

\begin{quote}
The girl sang a beautiful tune.
\end{quote}

Plural: \ar{ألحان}, \emph{alḥān}.

In grammar \ar{لحن} also means a mistake in speech.

\begin{cousinweb}
\textbf{\ar{لحن}} (\emph{laḥn}) is built on the native Arabic root \emph{l-ḥ-n}, \textquotedblleft{}to speak with a tone; to make a mistake in speech\textquotedblright{} (Lane). Lane gives laḥn both as a tone of voice and as a fault in speaking.
\end{cousinweb}

\begin{grammarlens}[title={What you've built}]
One more word for describing what you lived and felt.
\end{grammarlens}

\subsection*{Your turn}
\begin{itemize}
\raggedright
  \item \emph{Say it:} \emph{laḥn}
  \item \emph{Say it:} \emph{laḥn}, once more
  \item \emph{Say it:} \emph{ghannat al-bint laḥnan jamīlan.} — The girl sang a beautiful tune.
  \item \emph{Recall:} say the Arabic for to applaud, then the Arabic for acting, then say \emph{laḥn} again
\end{itemize}

\begin{tcolorbox}[breakable,colback=teal!4,colframe=teal!35!black,title={Before you move on}]
What does \ar{لحن} mean? (\textquotedblleft{}A melody, a tune\textquotedblright{}.) Say it once more.
\end{tcolorbox}
\section[\texorpdfstring{\ar{إيقاع} (īqāʿ)}{إيقاع (īqāʿ)}]{\ar{إيقاع} (īqāʿ) — a rhythm, a beat}

\label{lesson:AR-C543-iqa}



\begin{quote}
Before the new one: say the Arabic for acting, then the Arabic for a melody.
\end{quote}

\subsection*{You'll want to know: \ar{إيقاع}}
\textbf{\ar{إيقاع}} — \emph{īqāʿ} — \textquotedblleft{}a rhythm, a beat\textquotedblright{}.

\begin{quote}
\ar{رقص} \ar{الأطفال} \ar{على} \ar{إيقاع} \ar{الموسيقى}.
\end{quote}

\begin{quote}
\emph{raqaṣa al-aṭfāl ʿalā īqāʿ al-mūsīqā.}
\end{quote}

\begin{quote}
The children danced to the rhythm of the music.
\end{quote}

\begin{cousinweb}
\textbf{\ar{إيقاع}} (\emph{īqāʿ}) is built on the native Arabic root \emph{w-q-ʿ}, \textquotedblleft{}to fall, to happen\textquotedblright{} (Lane). Īqāʿ is the verbal noun of Form IV: making the beats fall.
\end{cousinweb}

\begin{grammarlens}[title={What you've built}]
One more word for describing what you lived and felt.
\end{grammarlens}

\subsection*{Your turn}
\begin{itemize}
\raggedright
  \item \emph{Say it:} \emph{īqāʿ}
  \item \emph{Say it:} \emph{īqāʿ}, once more
  \item \emph{Say it:} \emph{raqaṣa al-aṭfāl ʿalā īqāʿ al-mūsīqā.} — The children danced to the rhythm of the music.
  \item \emph{Recall:} say the Arabic for acting, then the Arabic for a melody, then say \emph{īqāʿ} again
\end{itemize}

\begin{tcolorbox}[breakable,colback=teal!4,colframe=teal!35!black,title={Before you move on}]
What does \ar{إيقاع} mean? (\textquotedblleft{}A rhythm, a beat\textquotedblright{}.) Say it once more.
\end{tcolorbox}
\section[\texorpdfstring{\ar{عزف} (ʿazafa)}{عزف (ʿazafa)}]{\ar{عزف} (ʿazafa) — to play an instrument}

\label{lesson:AR-C543-azafa}



\begin{quote}
Before the new one: say the Arabic for a melody, then the Arabic for a rhythm.
\end{quote}

\subsection*{You'll want to know: \ar{عزف}}
\textbf{\ar{عزف}} — \emph{ʿazafa} — \textquotedblleft{}to play an instrument\textquotedblright{}.

\begin{quote}
\ar{عزفت} \ar{أختي} \ar{أغنية} \ar{قديمة}.
\end{quote}

\begin{quote}
\emph{ʿazafat ukhtī ughniya qadīma.}
\end{quote}

\begin{quote}
My sister played an old song.
\end{quote}

Present: \ar{يعزف}, \emph{yaʿzif}.

It takes \ar{على} for the instrument.

\begin{cousinweb}
\textbf{\ar{عزف}} (\emph{ʿazafa}) is built on the native Arabic root \emph{ʿ-z-f}, \textquotedblleft{}to turn away; to make a sound\textquotedblright{} (Lane). Lane gives yaʿzifu as its present.
\end{cousinweb}

\begin{grammarlens}[title={What you've built}]
One more word for describing what you lived and felt.
\end{grammarlens}

\subsection*{Your turn}
\begin{itemize}
\raggedright
  \item \emph{Say it:} \emph{ʿazafa}
  \item \emph{Say it:} \emph{ʿazafa}, once more
  \item \emph{Say it:} \emph{ʿazafat ukhtī ughniya qadīma.} — My sister played an old song.
  \item \emph{Recall:} say the Arabic for a melody, then the Arabic for a rhythm, then say \emph{ʿazafa} again
\end{itemize}

\begin{tcolorbox}[breakable,colback=teal!4,colframe=teal!35!black,title={Before you move on}]
What does \ar{عزف} mean? (\textquotedblleft{}To play an instrument\textquotedblright{}.) Say it once more.
\end{tcolorbox}
\section[\texorpdfstring{\ar{ملحن} (mulaḥḥin)}{ملحن (mulaḥḥin)}]{\ar{ملحن} (mulaḥḥin) — a composer}

\label{lesson:AR-C543-mulahhin}



\begin{quote}
Before the new one: say the Arabic for a rhythm, then the Arabic for to play an instrument.
\end{quote}

\subsection*{You'll want to know: \ar{ملحن}}
\textbf{\ar{ملحن}} — \emph{mulaḥḥin} — \textquotedblleft{}a composer\textquotedblright{}.

\begin{quote}
\ar{كتب} \ar{الملحن} \ar{موسيقى} \ar{الفيلم}.
\end{quote}

\begin{quote}
\emph{kataba al-mulaḥḥin mūsīqā al-fīlm.}
\end{quote}

\begin{quote}
The composer wrote the music for the film.
\end{quote}

\begin{cousinweb}
\textbf{\ar{ملحن}} (\emph{mulaḥḥin}) is built on the native Arabic root \emph{l-ḥ-n}, \textquotedblleft{}to speak with a tone; to make a mistake in speech\textquotedblright{} (Lane). Mulaḥḥin is the participle of Form II, to set to a tune.
\end{cousinweb}

\begin{grammarlens}[title={What you've built}]
One more word for describing what you lived and felt.
\end{grammarlens}

\subsection*{Your turn}
\begin{itemize}
\raggedright
  \item \emph{Say it:} \emph{mulaḥḥin}
  \item \emph{Say it:} \emph{mulaḥḥin}, once more
  \item \emph{Say it:} \emph{kataba al-mulaḥḥin mūsīqā al-fīlm.} — The composer wrote the music for the film.
  \item \emph{Recall:} say the Arabic for a rhythm, then the Arabic for to play an instrument, then say \emph{mulaḥḥin} again
\end{itemize}

\begin{tcolorbox}[breakable,colback=teal!4,colframe=teal!35!black,title={Before you move on}]
What does \ar{ملحن} mean? (\textquotedblleft{}A composer\textquotedblright{}.) Say it once more.
\end{tcolorbox}
\section[\texorpdfstring{\ar{جوقة} (jawqa)}{جوقة (jawqa)}]{\ar{جوقة} (jawqa) — a choir}

\label{lesson:AR-C543-jawqa}



\begin{quote}
Before the new one: say the Arabic for to play an instrument, then the Arabic for a composer.
\end{quote}

\subsection*{You'll want to know: \ar{جوقة}}
\textbf{\ar{جوقة}} — \emph{jawqa} — \textquotedblleft{}a choir\textquotedblright{}.

\begin{quote}
\ar{تغني} \ar{ابنتي} \ar{في} \ar{جوقة} \ar{المدرسة}.
\end{quote}

\begin{quote}
\emph{tughannī ibnatī fī jawqat al-madrasa.}
\end{quote}

\begin{quote}
My daughter sings in the school choir.
\end{quote}

\begin{cousinweb}
\textbf{\ar{جوقة}} (\emph{jawqa}) is a word for a band or company of people, and today a choir; the origin is not settled.
\end{cousinweb}

\begin{grammarlens}[title={What you've built}]
One more word for describing what you lived and felt.
\end{grammarlens}

\subsection*{Your turn}
\begin{itemize}
\raggedright
  \item \emph{Say it:} \emph{jawqa}
  \item \emph{Say it:} \emph{jawqa}, once more
  \item \emph{Say it:} \emph{tughannī ibnatī fī jawqat al-madrasa.} — My daughter sings in the school choir.
  \item \emph{Recall:} say the Arabic for to play an instrument, then the Arabic for a composer, then say \emph{jawqa} again
\end{itemize}

\begin{tcolorbox}[breakable,colback=teal!4,colframe=teal!35!black,title={Before you move on}]
What does \ar{جوقة} mean? (\textquotedblleft{}A choir\textquotedblright{}.) Say it once more.
\end{tcolorbox}
\section[(dialogue)]{Words from chapters 536–543 — a review}

\label{lesson:AR-R543-review-536-543}



\begin{quote}
Say the Arabic for hunting and a fishing rod.
\end{quote}

\subsection*{Your turn}
Say these aloud:

\begin{itemize}
\raggedright
  \item hunting, a fishing rod, to go fishing, to unwind, to enjoy
  \item camping, diving, photography, pottery, a crossword
  \item embroidery, to embroider, knitting, to knit, to sew
  \item dice, a swing, to have fun, gardening, to swing
  \item reading, to subscribe, to meet up, passion, fondness
  \item an artist, an exhibition, a sculptor, to carve, a painter
  \item a stage play, a scene, applause, to applaud, acting
  \item a melody, a rhythm, to play an instrument, a composer, a choir
\end{itemize}

\begin{tcolorbox}[breakable,colback=teal!4,colframe=teal!35!black,title={Before you move on}]
Hunting? (\textbf{\ar{صيد}}, \emph{ṣayd}.) A fishing rod? (\textbf{\ar{صنارة}}, \emph{ṣinnāra}.) To go fishing? (\textbf{\ar{اصطاد}}, \emph{iṣṭāda}.) To unwind? (\textbf{\ar{استجم}}, \emph{istajamma}.) To enjoy? (\textbf{\ar{استمتع}}, \emph{istamtaʿa}.) Camping? (\textbf{\ar{تخييم}}, \emph{takhyīm}.) Diving? (\textbf{\ar{غوص}}, \emph{ghawṣ}.) Photography? (\textbf{\ar{تصوير}}, \emph{taṣwīr}.) Pottery? (\textbf{\ar{فخار}}, \emph{fakhkhār}.) A crossword? (\textbf{\ar{كلمات} \ar{متقاطعة}}, \emph{kalimāt mutaqāṭiʿa}.) Embroidery? (\textbf{\ar{تطريز}}, \emph{taṭrīz}.) To embroider? (\textbf{\ar{طرز}}, \emph{ṭarraza}.) Knitting? (\textbf{\ar{حياكة}}, \emph{ḥiyāka}.) To knit? (\textbf{\ar{حاك}}, \emph{ḥāka}.) To sew? (\textbf{\ar{خاط}}, \emph{khāṭa}.) Dice? (\textbf{\ar{نرد}}, \emph{nard}.) A swing? (\textbf{\ar{أرجوحة}}, \emph{urjūḥa}.) To have fun? (\textbf{\ar{تسلى}}, \emph{tasallā}.) Gardening? (\textbf{\ar{بستنة}}, \emph{bastana}.) To swing? (\textbf{\ar{تأرجح}}, \emph{taʾarjaḥa}.) Reading? (\textbf{\ar{قراءة}}, \emph{qirāʾa}.) To subscribe? (\textbf{\ar{اشترك}}, \emph{ishtaraka}.) To meet up? (\textbf{\ar{التقى}}, \emph{iltaqā}.) Passion? (\textbf{\ar{شغف}}, \emph{shaghaf}.) Fondness? (\textbf{\ar{ولع}}, \emph{walaʿ}.) An artist? (\textbf{\ar{فنان}}, \emph{fannān}.) An exhibition? (\textbf{\ar{معرض}}, \emph{maʿriḍ}.) A sculptor? (\textbf{\ar{نحات}}, \emph{naḥḥāt}.) To carve? (\textbf{\ar{نحت}}, \emph{naḥata}.) A painter? (\textbf{\ar{رسام}}, \emph{rassām}.) A stage play? (\textbf{\ar{مسرحية}}, \emph{masraḥiyya}.) A scene? (\textbf{\ar{مشهد}}, \emph{mashhad}.) Applause? (\textbf{\ar{تصفيق}}, \emph{taṣfīq}.) To applaud? (\textbf{\ar{صفق}}, \emph{ṣaffaqa}.) Acting? (\textbf{\ar{تمثيل}}, \emph{tamthīl}.) A melody? (\textbf{\ar{لحن}}, \emph{laḥn}.) A rhythm? (\textbf{\ar{إيقاع}}, \emph{īqāʿ}.) To play an instrument? (\textbf{\ar{عزف}}, \emph{ʿazafa}.) A composer? (\textbf{\ar{ملحن}}, \emph{mulaḥḥin}.) A choir? (\textbf{\ar{جوقة}}, \emph{jawqa}.)
\end{tcolorbox}
